% Generated draft: requires mathematical and source review.
\documentclass[11pt]{article}
\usepackage[margin=25mm]{geometry}
\usepackage{fontspec}
\setmainfont{Latin Modern Roman}
\usepackage{amsmath,amssymb,mathtools,unicode-math}
\setmathfont{Latin Modern Math}
\usepackage{enumitem,needspace,xurl,hyperref}
\setlength{\parindent}{0pt}
\setlength{\parskip}{5pt}
\setlength{\emergencystretch}{3em}
\newcommand{\N}{\mathbb{N}}
\newcommand{\Z}{\mathbb{Z}}
\newcommand{\Q}{\mathbb{Q}}
\newcommand{\R}{\mathbb{R}}
\newcommand{\C}{\mathbb{C}}
\newcommand{\F}{\mathbb{F}}
\newcommand{\A}{\mathbb{A}}
\newcommand{\PP}{\mathbb{P}}
\newcommand{\E}{\mathbb{E}}
\newcommand{\eps}{\varepsilon}
\newcommand{\dd}{\,\mathrm{d}}
\begin{document}
\section*{UnsolvedMath ID: 30003701}
\textbf{Source problem code:} catalogue entry
\subsection{Definitions and mathematical statement}
The supplied record does not specify a base field or its characteristic. The following formulation makes that missing parameter explicit; it does not assert that every base field is intended.

Fix a field \(k\) and a positive integer \(n\). Let
\[
\mathcal P_{n,k}=\mathbb A_k^n,
\qquad
p_a(T)=T^n+a_1T^{n-1}+\cdots+a_n,
\]
and define the incidence variety
\[
\widetilde{\mathcal P}_{n,k}
=\{(a,t)\in\mathbb A_k^n\times\mathbb A_k^1:p_a(t)=0\}.
\]
The projection \(\pi_{n,k}(a,t)=a\) is the finite morphism recording a monic polynomial and forgetting its chosen root. It includes polynomials with repeated roots; its restriction to the locus of nonzero discriminant is finite étale. The resolvent degree below concerns the generic polynomial, so it imposes no requirement that a root-selection construction extend over the discriminant locus.

Here is an algebraic field-tower definition of the invariant, conditional on the choice of \(k\). Put
\[
K=k(a_1,\ldots,a_n),\qquad L=K(\alpha),\qquad p_a(\alpha)=0,
\]
where the coefficients are algebraically independent over \(k\) and \(\alpha\) is a root of the generic polynomial. Define \(\operatorname{RD}_k(L/K)\) to be the least nonnegative integer \(d\) for which there is a tower of finite field extensions
\[
K=K_0\subset K_1\subset\cdots\subset K_r
\]
and a \(K\)-embedding \(L\hookrightarrow K_r\), with the following property: for every \(i\), there are finitely generated field extensions \(F_i/k\) and \(E_i/k\), with \(E_i/F_i\) finite and \(\operatorname{trdeg}_k F_i\le d\), and compatible \(k\)-embeddings \(F_i\hookrightarrow K_{i-1}\) and \(E_i\hookrightarrow K_i\), such that \(K_i=K_{i-1}E_i\). Set
\[
\operatorname{RD}_k(\widetilde{\mathcal P}_{n,k}/\mathcal P_{n,k})
=\operatorname{RD}_k(L/K).
\]

The requested problem is to determine this resolvent degree as a function of \(n\), once the intended base field and conventions have been fixed. The record additionally asks to decide the conjectures it calls Hilbert's sextic, thirteenth, and octic conjectures. It supplies neither their numerical target values nor their precise formulations, so no particular equalities or inequalities for those named conjectures can be stated from this record alone. The base-field qualification and the missing conjecture formulations remain unresolved.
\subsection{Short English statement}
Determine how many independent parameters are needed in successive algebraic auxiliary problems to obtain a root of a generic monic polynomial of degree \(n\). The source also requests decisions on Hilbert's sextic, thirteenth, and octic conjectures, but leaves the base field and the precise conjectural assertions unspecified.
\subsection{Sources}
\begin{itemize}
\item Local UnsolvedMath registry, permanent ID \(30003701\), title \emph{Resolvent Degree of Polynomial Root Covers}; exact source problem code \texttt{catalogue entry}. The supplied collection is \emph{Oberwolfach Reports Open Problems}. No individual proposer is identified; the supplied proposed year is 2018.
\item Supplied report citation: \emph{Topology of Arrangements and Representation Stability (2019)}, Oberwolfach Reports, DOI \texttt{10.4171/owr/2018/2}, pp.\ 43--123. \url{https://doi.org/10.4171/owr/2018/2}. The record describes these reports as documenting mathematical workshops and their open-problem discussions. The year in the supplied report title differs from the supplied proposed year and the year appearing in the DOI.
\item Supplied literature reference: Edens and Reichstein, \emph{Hilbert thirteen in prime characteristic}. \url{https://arxiv.org/abs/2406.15954}.
\item Supplied literature reference: Farb and Wolfson, \emph{resolvent degree}. \url{https://arxiv.org/abs/1803.04063}.
\item Supplied literature-check date: \texttt{2026-08-21}. The supplied dated curation labels the general resolvent-degree problem and the named sextic, thirteenth, and octic conjectures open, and reports that only upper bounds are known for the target values it references. These are supplied literature assessments, not independently verified conclusions. The background explicitly requests independent verification of time-sensitive status and cited claims.
\item The registry labels the statement \texttt{corrected\_verified}, explaining that its mathematical content was preserved while notation, typography, or omitted definitions were repaired. This is supplied registry metadata, not mathematical verification of this draft.
\end{itemize}
\end{document}
